\documentclass{beamer}

% Theme choice
\usetheme{Madrid} % You can change to e.g., Warsaw, Berlin, CambridgeUS, etc.

% Encoding and font
\usepackage[utf8]{inputenc}
\usepackage[T1]{fontenc}

% Graphics and tables
\usepackage{graphicx}
\usepackage{booktabs}

% Code listings
\usepackage{listings}
\lstset{
basicstyle=\ttfamily\small,
keywordstyle=\color{blue},
commentstyle=\color{gray},
stringstyle=\color{red},
breaklines=true,
frame=single
}

% Math packages
\usepackage{amsmath}
\usepackage{amssymb}

% Colors
\usepackage{xcolor}

% TikZ and PGFPlots
\usepackage{tikz}
\usepackage{pgfplots}
\pgfplotsset{compat=1.18}
\usetikzlibrary{positioning}

% Hyperlinks
\usepackage{hyperref}

% Title information
\title{Week 8: Automated Testing Tools}
\author{Your Name}
\institute{Your Institution}
\date{\today}

\begin{document}

\frame{\titlepage}

\begin{frame}[fragile]
    \frametitle{Introduction to Automated Testing Tools}
    \begin{block}{Overview of Automated Testing}
        Automated testing is a pivotal component in modern software development that facilitates 
        the efficient testing of applications, ensuring reliability, performance, and quality.
    \end{block}
\end{frame}

\begin{frame}[fragile]
    \frametitle{Importance of Automated Testing}
    \begin{enumerate}
        \item \textbf{Increased Efficiency}: 
        \begin{itemize}
            \item Automated tests execute significantly faster than manual tests.
            \item Example: A full suite of automated tests takes minutes, whereas manual tests could take hours.
        \end{itemize}
        
        \item \textbf{Consistency and Reliability}: 
        \begin{itemize}
            \item Reduces human error and ensures tests are executed similarly each time.
            \item Consistent testing environments prevent discrepancies.
        \end{itemize}
        
        \item \textbf{Cost-Effective in the Long Run}:
        \begin{itemize}
            \item Initial costs may be high, but long-term savings from reduced testing time are significant.
        \end{itemize}
        
        \item \textbf{Enhanced Test Coverage}: 
        \begin{itemize}
            \item More test cases and scenarios can be run compared to manual testing.
            \item Supports regression testing to check new code changes against existing functionality.
        \end{itemize}
        
        \item \textbf{Supports CI/CD}: 
        \begin{itemize}
            \item Integrates seamlessly into Continuous Integration/Continuous Deployment pipelines.
        \end{itemize}
    \end{enumerate}
\end{frame}

\begin{frame}[fragile]
    \frametitle{Examples of Automated Testing Tools}
    \begin{itemize}
        \item \textbf{Selenium}: 
        Used primarily for web application testing, supports multiple programming languages.
        
        \item \textbf{JUnit}: 
        A popular framework for unit testing in Java.
        
        \item \textbf{TestNG}: 
        Inspired by JUnit, it allows for more flexible test configurations and data-driven tests.
        
        \item \textbf{Cypress}: 
        A modern tool for end-to-end testing specifically designed for web applications.
    \end{itemize}
\end{frame}

\begin{frame}[fragile]
    \frametitle{Conclusion}
    Automated testing is essential in today's fast-paced software development environment. 
    It reduces manual effort, increases accuracy, and enables rapid feedback loops, significantly 
    enhancing software quality and helping deliver products that meet user expectations.
\end{frame}

\begin{frame}[fragile]
    \frametitle{What is Unit Testing? - Definition}
    \begin{block}{Definition}
        Unit testing is a software testing method where individual components (or "units") of a program are tested in isolation from the rest of the application. The primary goal is to validate that each unit of the software performs as designed.
    \end{block}
\end{frame}

\begin{frame}[fragile]
    \frametitle{What is Unit Testing? - Objectives}
    \begin{itemize}
        \item \textbf{Verification}: Ensure that each component works correctly and meets specifications.
        \item \textbf{Early Bug Detection}: Identify issues early, reducing the propagation of defects.
        \item \textbf{Simplifies Integration}: Makes integration into larger systems easier.
        \item \textbf{Facilitates Code Change}: Provides safety for code changes or refactoring.
        \item \textbf{Documentation}: Serves as documentation of expected unit behavior.
    \end{itemize}
\end{frame}

\begin{frame}[fragile]
    \frametitle{What is Unit Testing? - Key Points and Example}
    \begin{block}{Key Points}
        \begin{itemize}
            \item \textbf{Granularity}: Each test focuses on a single piece or function of the codebase.
            \item \textbf{Automation}: Can be automated with various testing tools.
            \item \textbf{Framework Support}: Supported by many programming languages (e.g., JUnit, NUnit, PyTest).
        \end{itemize}
    \end{block}
    
    \begin{block}{Example}
        \begin{lstlisting}[language=python]
import unittest

def add(a, b):
    return a + b

class TestAddFunction(unittest.TestCase):
    
    def test_add_positive_numbers(self):
        self.assertEqual(add(2, 3), 5)

    def test_add_negative_numbers(self):
        self.assertEqual(add(-1, -1), -2)

if __name__ == '__main__':
    unittest.main()
        \end{lstlisting}
    \end{block}
\end{frame}

\begin{frame}[fragile]
    \frametitle{Introduction to JUnit - Overview}
    \begin{block}{What is JUnit?}
        JUnit is a widely used testing framework in Java that supports the creation and execution of unit tests,
        allowing developers to write tests for small chunks of code (units) to ensure that they work as expected.
    \end{block}
\end{frame}

\begin{frame}[fragile]
    \frametitle{Introduction to JUnit - History}
    \begin{itemize}
        \item \textbf{First Release:} Created by Kent Beck and Erich Gamma in 1997 as part of Test-Driven Development (TDD).
        \item \textbf{Evolution:} Significant versions include JUnit 4 and JUnit 5; JUnit 5 introduced a modernized architecture.
    \end{itemize}
\end{frame}

\begin{frame}[fragile]
    \frametitle{Introduction to JUnit - Role in Java Unit Testing}
    \begin{itemize}
        \item \textbf{Promotes Good Coding Habits:} Encourages writing testable code and improves code quality.
        \item \textbf{Automation:} Facilitates automated testing, enabling seamless test execution during builds.
        \item \textbf{Integration:} Easily integrates with build tools like Maven and Gradle, as well as IDEs such as IntelliJ IDEA and Eclipse.
    \end{itemize}
\end{frame}

\begin{frame}[fragile]
    \frametitle{Introduction to JUnit - Structure}
    \begin{enumerate}
        \item \textbf{Test Classes:} Contain methods that perform tests; usually correspond to the application class being tested.
        \item \textbf{Test Methods:} Decorated with the \texttt{@Test} annotation, executing assertions to verify code behavior.
    \end{enumerate}
\end{frame}

\begin{frame}[fragile]
    \frametitle{Introduction to JUnit - Core Features}
    \begin{itemize}
        \item \textbf{Annotations:} Key annotations include \texttt{@Test}, \texttt{@Before}, \texttt{@After}, \texttt{@BeforeClass}, and \texttt{@AfterClass}.
        \item \textbf{Assertions:} Used to assert expected results:
            \begin{lstlisting}
            assertEquals(expectedValue, actualValue);
            assertTrue(condition);
            assertFalse(condition);
            \end{lstlisting}
        \item \textbf{Test Suites:} Groups multiple test classes together for collective testing.
    \end{itemize}
\end{frame}

\begin{frame}[fragile]
    \frametitle{Introduction to JUnit - Example Code Snippet}
    \begin{lstlisting}[language=Java]
import org.junit.Assert;
import org.junit.Test;

public class CalculatorTest {
    @Test
    public void testAdd() {
        Calculator calc = new Calculator();
        Assert.assertEquals(5, calc.add(2, 3));
    }
}
    \end{lstlisting}
    \begin{block}{Explanation}
        This example verifies that the \texttt{add} method of a \texttt{Calculator} class correctly sums two numbers.
        Assertions check if the expected outcome matches the actual output.
    \end{block}
\end{frame}

\begin{frame}[fragile]
    \frametitle{Writing and Running JUnit Tests - Overview}
    JUnit is a powerful testing framework for Java that allows developers to write and run repeatable tests efficiently. 
    This slide covers the step-by-step process of creating unit tests using JUnit, focusing on essential annotations and assertions.
\end{frame}

\begin{frame}[fragile]
    \frametitle{Writing JUnit Tests - Step-by-Step Process}
    \begin{enumerate}
        \item \textbf{Setup Your JUnit Environment}
        \begin{itemize}
            \item Ensure your IDE has JUnit included; if not, add the JUnit library (e.g., JUnit 5).
        \end{itemize}
        
        \item \textbf{Create a Test Class}
        \begin{itemize}
            \item For every class to test, create a corresponding test class (e.g., if `Calculator`, then `CalculatorTest`).
        \end{itemize}
        
        \begin{lstlisting}[language=Java]
import static org.junit.jupiter.api.Assertions.*;
import org.junit.jupiter.api.Test;

public class CalculatorTest {
    // Test methods will go here
}
        \end{lstlisting}
    \end{enumerate}
\end{frame}

\begin{frame}[fragile]
    \frametitle{Annotations and Assertions in JUnit}
    \begin{enumerate}
        \setcounter{enumi}{2}
        \item \textbf{Use the @Test Annotation}
        \begin{lstlisting}[language=Java]
@Test
void testAddition() {
    // Your test logic will go here
}
        \end{lstlisting}

        \item \textbf{Write Assertions}
        Use assertions to validate your test results. Common assertions include:
        \begin{itemize}
            \item \texttt{assertEquals(expected, actual)} - checks if two values are equal.
            \item \texttt{assertTrue(condition)} - verifies that a condition is true.
            \item \texttt{assertFalse(condition)} - verifies that a condition is false.
            \item \texttt{assertNull(object)} - checks if an object is null.
            \item \texttt{assertNotNull(object)} - checks if an object is not null.
        \end{itemize}

        \textbf{Example of a Test Method with Assertions:}
        \begin{lstlisting}[language=Java]
@Test
void testAddition() {
    Calculator calc = new Calculator();
    int result = calc.add(2, 3);
    assertEquals(5, result); // Should pass
}
        \end{lstlisting}
    \end{enumerate}
\end{frame}

\begin{frame}[fragile]
    \frametitle{Introduction to pytest - Overview}
    \begin{block}{Overview of pytest}
        pytest is a robust testing framework for Python that simplifies the process of writing and executing test cases. 
        It is widely adopted due to its flexibility, scalability, and user-friendly features, making it a preferred choice for both beginners and advanced Python developers.
    \end{block}
\end{frame}

\begin{frame}[fragile]
    \frametitle{Introduction to pytest - Uniqueness}
    \begin{itemize}
        \item \textbf{Simple Syntax:} Write tests using simple assert statements.
        \item \textbf{Support for Fixtures:} Enables reusable setups, promoting DRY principles.
        \item \textbf{Rich Plugin Architecture:} Easily extended with various plugins for diverse testing needs.
        \item \textbf{Automatic Test Discovery:} Automatically finds tests via naming conventions.
    \end{itemize}
\end{frame}

\begin{frame}[fragile]
    \frametitle{Introduction to pytest - Advantages}
    \begin{itemize}
        \item \textbf{Flexibility:} Supports both simple unit tests and complex functional test suites.
        \item \textbf{Descriptive Output:} Provides detailed test result outputs for easier debugging.
        \item \textbf{Parameterized Testing:} Run the same test with different parameters.
        \item \textbf{Compatibility:} Works seamlessly with existing unittest and nose tests.
    \end{itemize}
\end{frame}

\begin{frame}[fragile]
    \frametitle{Introduction to pytest - Example Test}
    \begin{block}{Example of a Simple pytest Test}
        \begin{lstlisting}[language=Python]
def add(x, y):
    return x + y

def test_add():
    assert add(3, 2) == 5
    assert add(-1, 1) == 0
        \end{lstlisting}
    \end{block}
\end{frame}

\begin{frame}[fragile]
    \frametitle{Introduction to pytest - Conclusion}
    \begin{block}{Conclusion}
        pytest stands out as a leading testing framework in Python due to its simplicity, versatility, and rich feature set.
        Familiarizing yourself with pytest can drastically improve the efficiency of your testing processes.
    \end{block}
\end{frame}

\begin{frame}[fragile]
    \frametitle{Writing and Running pytest Tests - Introduction}
    Pytest is a powerful testing framework in Python that allows you to write simple and scalable test cases.
    \begin{itemize}
        \item Encourages the use of fixtures and assertions.
        \item Provides extensive capabilities for unit testing.
    \end{itemize}
\end{frame}

\begin{frame}[fragile]
    \frametitle{Creating Unit Tests}
    Unit tests in pytest are defined in files that start with \texttt{test\_} or end with \texttt{\_test.py}. Each test function should also begin with \texttt{test\_}.
    
    \begin{block}{Example: Basic Unit Test}
    \begin{lstlisting}[language=Python]
# test_calculator.py
def add(a, b):
    return a + b

def test_add():
    assert add(2, 3) == 5  # This test should pass
    assert add(-1, 1) == 0  # This test should pass
    assert add(2, -2) == 0  # This test should pass
    \end{lstlisting}
    \end{block}
    
    - In this example, \texttt{test\_add} verifies the behavior of the \texttt{add} function.
\end{frame}

\begin{frame}[fragile]
    \frametitle{Running Tests and Using Fixtures}
    To run tests, execute the command:
    \begin{lstlisting}[language=bash]
pytest
    \end{lstlisting}
    \begin{itemize}
        \item Automatically discovers all test files.
        \item Shows summary of passed and failed tests.
    \end{itemize}

    \textbf{Common Command Options:}
    \begin{itemize}
        \item \texttt{-v}: Verbose mode for detailed output.
        \item \texttt{-k <expression>}: Run tests that match the expression.
    \end{itemize}
    
    \begin{block}{Example: Using Fixtures}
    \begin{lstlisting}[language=Python]
import pytest

@pytest.fixture
def sample_data():
    return [1, 2, 3]

def test_length(sample_data):
    assert len(sample_data) == 3  # This test should pass
    \end{lstlisting}
    \end{block}
\end{frame}

\begin{frame}[fragile]
    \frametitle{Comparing JUnit and pytest - Introduction}
    \begin{itemize}
        \item JUnit and pytest are popular testing frameworks for unit testing in Java and Python.
        \item Understanding their similarities and differences is essential for selecting the right tool based on programming language and project requirements.
    \end{itemize}
\end{frame}

\begin{frame}[fragile]
    \frametitle{Comparing JUnit and pytest - Similarities}
    \begin{itemize}
        \item Both frameworks provide strong support for unit testing and assertions.
        \item Support test organization through test suites: 
        \begin{itemize}
            \item JUnit uses \texttt{@Suite}
            \item pytest uses \texttt{@pytest.mark.parametrize}
        \end{itemize}
        \item Enable running tests automatically, with integration for Continuous Integration (CI) tools.
    \end{itemize}
\end{frame}

\begin{frame}[fragile]
    \frametitle{Comparing JUnit and pytest - Differences}
    \textbf{Language and Syntax}
    \begin{itemize}
        \item \textbf{JUnit (Java)}:
        \begin{lstlisting}[language=java]
import org.junit.Test;
import static org.junit.Assert.assertEquals;

public class CalculatorTest {
    @Test
    public void testAdd() {
        Calculator calculator = new Calculator();
        assertEquals(5, calculator.add(2, 3));
    }
}
\end{lstlisting}
        \item \textbf{pytest (Python)}:
        \begin{lstlisting}[language=python]
import pytest

def test_add():
    assert add(2, 3) == 5
\end{lstlisting}
    \end{itemize}
\end{frame}

\begin{frame}[fragile]
    \frametitle{Comparing JUnit and pytest - Features}
    \begin{itemize}
        \item \textbf{JUnit}:
        \begin{itemize}
            \item Emphasizes structured test cases, requires \texttt{assert} methods.
            \item More verbose setup with test suites and runner configurations.
        \end{itemize}
        \item \textbf{pytest}:
        \begin{itemize}
            \item Supports fixtures for reusable setup/teardown code.
            \item Offers rich plugins for mocking, coverage, and parallel test execution.
        \end{itemize}
    \end{itemize}
\end{frame}

\begin{frame}[fragile]
    \frametitle{Comparing JUnit and pytest - Aesthetics & Conclusion}
    \begin{itemize}
        \item \textbf{JUnit}: More complex, ideal for developers with Java background.
        \item \textbf{pytest}: Simpler and more flexible, beginner-friendly for Python developers.
    \end{itemize}
    
    \textbf{Key Points:}
    \begin{enumerate}
        \item Choose based on Language: JUnit for Java, pytest for Python.
        \item Consider Syntax Preference: pytest's function-based approach may appeal for its simplicity.
        \item Leverage Features: Use pytest fixtures/plugins for more robust testing scenarios.
    \end{enumerate}
\end{frame}

\begin{frame}[fragile]
    \frametitle{Best Practices in Automated Testing - Introduction}
    \begin{itemize}
        \item Automated testing is essential for maintaining software quality.
        \item Best practices ensure that automated tests are:
        \begin{itemize}
            \item Effective
            \item Maintainable
            \item Efficient
        \end{itemize}
    \end{itemize}
\end{frame}

\begin{frame}[fragile]
    \frametitle{Best Practices in Automated Testing - Organize Your Tests}
    \begin{block}{Structure Your Tests}
        \begin{itemize}
            \item \textbf{Folder Structure:} Separate tests by feature or module.
            \item \textbf{Example:}
            \begin{lstlisting}[language=Python]
/tests
  /unit
    test_module1.py
    test_module2.py
  /integration
    test_integration1.py
            \end{lstlisting}
            \item \textbf{Naming Conventions:} Descriptive names for test files/functions.
            \item \textbf{Example:} Use \texttt{test\_login\_user.py} instead of \texttt{test1.py}.
        \end{itemize}
    \end{block}
\end{frame}

\begin{frame}[fragile]
    \frametitle{Best Practices in Automated Testing - Maintainability}
    \begin{block}{Write Clear and Concise Tests}
        \begin{itemize}
            \item Ensure tests are easy to read and understand.
            \item \textbf{Example Assertion:}
            \begin{lstlisting}[language=Python]
assert user.is_authenticated() == True, "User should be authenticated"
            \end{lstlisting}
        \end{itemize}
        \begin{block}{Reduce Duplication}
            \begin{itemize}
                \item Utilize setup code and reusable functions.
                \item \textbf{Example:} Use fixtures in pytest for setup.
                \begin{lstlisting}[language=Python]
@pytest.fixture
def create_user():
    return User(username="testuser")
                \end{lstlisting}
            \end{itemize}
        \end{block}
    \end{block}
\end{frame}

\begin{frame}[fragile]
    \frametitle{Best Practices in Automated Testing - Continuous Integration}
    \begin{block}{Automate Test Execution}
        \begin{itemize}
            \item Integrate automated tests into your CI pipeline.
            \item \textbf{Example Tools:} Jenkins, GitHub Actions.
            \item Run tests automatically on code commits for immediate feedback.
        \end{itemize}
    \end{block}
\end{frame}

\begin{frame}[fragile]
    \frametitle{Best Practices in Automated Testing - Conclusion}
    \begin{itemize}
        \item Proper organization enhances readability and maintainability.
        \item Consistent naming conventions and structures improve navigation.
        \item Regular refactoring keeps the test suite relevant and efficient.
        \item Integrating with CI helps catch issues early in the development process.
    \end{itemize}
    \begin{block}{Key Takeaway}
        Implementing these best practices will lead to higher software quality and more dependable applications.
    \end{block}
\end{frame}

\begin{frame}[fragile]
    \frametitle{Integrating Tests into Continuous Integration (CI)}
    Continuous Integration (CI) is a development practice where developers frequently integrate code changes into a shared repository. Automated testing helps ensure that new changes do not break existing functionality.
    \begin{itemize}
        \item Integrate JUnit for Java projects.
        \item Integrate pytest for Python projects.
        \item Implement CI pipelines for improved efficiency.
    \end{itemize}
\end{frame}

\begin{frame}[fragile]
    \frametitle{Key Concepts in CI}
    \begin{enumerate}
        \item \textbf{What is CI?}
        \begin{itemize}
            \item Automates code integration and testing.
            \item Enhances collaboration and code quality.
            \item Tools include Jenkins, CircleCI, GitHub Actions.
        \end{itemize}
        
        \item \textbf{Why use Automated Testing in CI?}
        \begin{itemize}
            \item Immediate feedback after changes.
            \item Detect bugs early in the development cycle.
        \end{itemize}
    \end{enumerate}
\end{frame}

\begin{frame}[fragile]
    \frametitle{Integration Steps for JUnit (Java)}
    \textbf{Setting Up Your CI Pipeline}
    \begin{itemize}
        \item Choose a CI service (e.g., Jenkins, Travis CI).
        \item Create a configuration file in your project root.
    \end{itemize}
    
    \textbf{Example: JUnit in Jenkins}
    \begin{lstlisting}[language=groovy]
    pipeline {
        agent any
        stages {
            stage('Build') {
                steps {
                    sh 'mvn clean install'
                }
            }
            stage('Test') {
                steps {
                    sh 'mvn test'
                }
            }
        }
        post {
            always {
                junit 'target/surefire-reports/*.xml'
            }
        }
    }
    \end{lstlisting}
\end{frame}

\begin{frame}[fragile]
    \frametitle{Integration Steps for pytest (Python)}
    \textbf{Setting Up Your CI Pipeline for pytest}
    \begin{itemize}
        \item Ensure pytest is installed (use: pip install pytest).
        \item Create a configuration file in your project root.
    \end{itemize}

    \textbf{Example: pytest in GitHub Actions}
    \begin{lstlisting}[language=yaml]
    name: CI
    on: [push, pull_request]
    jobs:
      test:
        runs-on: ubuntu-latest
        steps:
          - name: Checkout code
            uses: actions/checkout@v2
          - name: Set up Python
            uses: actions/setup-python@v2
            with:
              python-version: '3.x'
          - name: Install dependencies
            run: |
              pip install -r requirements.txt
          - name: Run tests
            run: |
              pytest --junitxml=results.xml
          - name: Upload test results
            uses: actions/upload-artifact@v2
            with:
              name: Test Results
              path: results.xml
    \end{lstlisting}
\end{frame}

\begin{frame}[fragile]
    \frametitle{Best Practices in CI Testing}
    \begin{itemize}
        \item \textbf{Test Coverage:} Aim for high coverage but prioritize critical functionalities.
        \item \textbf{Failure Notifications:} Configure notifications to alert the team during test failures.
        \item \textbf{Keep Tests Fast:} Avoid bottlenecks in the CI process with quick tests.
    \end{itemize}
\end{frame}

\begin{frame}[fragile]
    \frametitle{Conclusion}
    Integrating JUnit and pytest testing frameworks into CI pipelines enhances software quality and improves development speed. Automation allows teams to focus on feature development while maintaining code integrity.
\end{frame}

\begin{frame}[fragile]
    \frametitle{Conclusion and Future Perspectives - Key Points Recap}
    \begin{enumerate}
        \item \textbf{Automated Testing Importance}:
        \begin{itemize}
            \item Enhances software quality by identifying bugs during development.
            \item Advantages: increased test coverage, faster feedback loops, reduced human error.
        \end{itemize}

        \item \textbf{Types of Automated Testing Tools}:
        \begin{itemize}
            \item \textbf{Unit Testing Frameworks}: JUnit (Java), pytest (Python).
            \item \textbf{Integration Testing Tools}: Postman, SoapUI.
            \item \textbf{End-to-End Testing}: Selenium, Cypress.
        \end{itemize}

        \item \textbf{Integration with CI/CD Pipelines}:
        \begin{itemize}
            \item Automated tests must be integrated into CI/CD pipelines to ensure code quality.
            \item Example: Using GitHub Actions to run pytest on code push.
        \end{itemize}
    \end{enumerate}
\end{frame}

\begin{frame}[fragile]
    \frametitle{Conclusion and Future Perspectives - Future Trends}
    \begin{enumerate}
        \item \textbf{AI and Machine Learning}:
        \begin{itemize}
            \item AI-driven tools can create, execute, and analyze tests automatically (e.g., Test.ai).
        \end{itemize}

        \item \textbf{Test Automation for Low-Code Platforms}:
        \begin{itemize}
            \item Specialized tools catering to low-code/no-code platforms emerging.
        \end{itemize}

        \item \textbf{Shift-Left Testing}:
        \begin{itemize}
            \item Encourages writing tests earlier in the development process for quicker defect identification.
        \end{itemize}
    \end{enumerate}
\end{frame}

\begin{frame}[fragile]
    \frametitle{Conclusion and Future Perspectives - Conclusion}
    \begin{block}{Conclusion}
        \begin{itemize}
            \item Evolution of automated testing tools boosts software quality and efficiency.
            \item Future advancements in AI, CI, and security redefine testing practices.
            \item Staying updated with trends is critical for competitive advantage in software development.
        \end{itemize}
    \end{block}
\end{frame}


\end{document}